\documentclass[10pt,a4paper]{amsart}
\usepackage[margin=19mm]{geometry}
\usepackage{amsmath,amssymb,mathtools}
\usepackage[scr=rsfs,cal=euler]{mathalfa}
\usepackage{xfrac}
\usepackage[unicode,hidelinks,bookmarks=false]{hyperref}
\newcommand{\ie}{i.e.\ }
\newcommand{\cf}{cf.\ }
\newcommand{\resp}{respectively\ }
\newcommand{\Subsection}{\subsection}
\providecommand{\nobreakdash}{\nobreak\hbox{-}}
\setlength{\emergencystretch}{2em}
\multlinegap0pt
\usepackage{bm}
\usepackage{bbm}
\usepackage{dsfont}
\usepackage{xspace}
\usepackage{cancel}
\usepackage{mathtools}
\usepackage{cleveref}
\usepackage{amsthm}
\makeatletter
\@ifundefined{theorem}{\newtheorem{theorem}{Theorem}}{}
\@ifundefined{lemma}{\newtheorem{lemma}[theorem]{Lemma}}{}
\makeatother
\def\BB{\mathbb{B}}
\def\CC{\mathbb{C}}
\def\FF{\mathbb{F}}
\def\NN{\mathbb{N}}
\def\a{\alpha}
\def\dju{\bigsqcup_{n=1}^{\infty}}
\def\nc{\mathrm{nc}}
\def\p{\varphi}
\def\s{\sigma}
\newcommand\ncm[1]{(\mathscr{A}_{#1}^{\dagger})_+}
\title[A Noncommutative Szeg\H{o}-Type Theorem on the Row-Ball]{A Noncommutative Szeg\H{o}-Type Theorem on the Row-Ball}
\author{Connor J. Gauntlett, David P. Kimsey}
\date{Source: arXiv:2507.19453v1 (CC BY 4.0)}
\begin{document}
\maketitle

\noindent\emph{Source and licence.} A Noncommutative Szeg\H{o}-Type Theorem on the Row-Ball. Version arXiv:2507.19453v1 (CC BY 4.0), distributed under the Creative Commons Attribution 4.0 International licence, as recorded in the arXiv licence metadata for this exact version and shown on the abstract page https://arxiv.org/abs/2507.19453v1. Full rights notice: licenses/arxiv-2507.19453-CC-BY-4.0.txt.\par\smallskip

\noindent\textit{Editorial scope.} Counted unit: the Lemma whose source label is \texttt{lem:ZerosBoundary} in the article (source0.tex lines 1603--1610, source label \texttt{lem:ZerosBoundary}), delivered together with the article's own complete original proof of it (source0.tex lines 1611--1662, one \texttt{proof} environment of 52 lines). No mathematics is written, repaired or re-derived: statement and proof are the article's own text line for line. Required context retained uncounted: lem:ZerosOutside (lines 1544--1547); lem:ZerosInside (lines 1578--1584). Every definition, display and label of those lines is retained unchanged. Omitted from this package and not redistributed: 1--1543, 1548--1577, 1585--1602, 1663--1696. Nothing inside a retained excerpt is omitted or altered: every retained body is delivered exactly as the article's lines are written, so no omission receipt of any kind accompanies this package. Citations: the retained excerpts carry no citation command at all, so no bibliography entry of the article is transcribed and no reference is redistributed. Reference adaptation: none was needed and none was made; every reference inside the delivered unit resolves inside it, so no unresolved reference or citation is produced. Printed slips kept literal: no printed word, symbol, hypothesis or number of the counted statement or of its proof was altered. Layout and class: the article's own class is \texttt{amsart} with packages including inputenc, mathtools, graphicx, xcolor, float, minipage-marginpar, blindtext, amssymb, amsmath, amsthm, cleveref, enumitem, cjhebrew, titletoc; the wrapper is the lane's pinned \texttt{amsart} editorial wrapper at 10pt on A4, which restates the theorem-like environments the retained text needs and adds the article's own macro definitions that the retained text needs (\texttt{\textbackslash BB}, \texttt{\textbackslash CC}, \texttt{\textbackslash FF}, \texttt{\textbackslash NN}, \texttt{\textbackslash a}, \texttt{\textbackslash dju}, \texttt{\textbackslash nc}, \texttt{\textbackslash ncm}, \texttt{\textbackslash p}, \texttt{\textbackslash s}), with extra wrapper packages (bm, bbm, dsfont, xspace, cancel, mathtools, cleveref). The delivered numbering is the unit's own. Assets: no retained span contains a figure environment, a graphics-inclusion command, a PDF-page inclusion, a table or a TikZ picture; no image file is copied into this package, the package declares no asset dependency and no image file is redistributed with it. Licence: version arXiv:2507.19453v1 of this article is distributed under the Creative Commons Attribution 4.0 International licence, as recorded in the arXiv licence metadata for this exact version and shown on the abstract page https://arxiv.org/abs/2507.19453v1. A full rights notice is delivered at licenses/arxiv-2507.19453-CC-BY-4.0.txt. Characters: every retained excerpt is pure ASCII, so the delivered engine, XeLaTeX with its default Latin Modern text font, reports no missing character.
\begin{lemma}
    \label{lem:ZerosOutside}
    Let \(\mu \in \ncm{d}\) be a non-trivial probability nc measure with associated orthonormal polynomials \((\p_\s)_{\s \in \FF_d^+}\). For any \(n \in \NN\), if \(Z \in \CC_\nc^d\) is such that \(\det \p_{\s(n)}^\#(Z^*) = 0\), then \(Z\) lies outside the row-ball.
\end{lemma}

\begin{lemma}
    \label{lem:ZerosInside}
    Let \(\mu \in \ncm{d}\) be a non-trivial probability nc measure with associated orthonormal polynomials \((\p_\s)_{\s \in \FF_d^+}\). For any \(n \in \NN\), if \(Z \in \CC_\nc^d\) is such that \(\det \sum_{\lvert \s \rvert = n} \p_\s(Z^*)^* \p_\s(Z^*) = 0\), then \(Z\) lies outside the set
    \[
        \mathrm{Ext}(\overline{\BB}^d_{\nc}) := \dju \left\{(Z_1, \ldots, Z_d) \in (\CC^{n \times n})^d : Z_1Z_1^* + \cdots + Z_dZ_d^* > I_n \right\}.
    \]
\end{lemma}

\begin{lemma}
    \label{lem:ZerosBoundary}
    Let \(\mu \in \ncm{d}\) be a non-trivial probability nc measure with associated orthonormal polynomials \((\p_\s)_{\s \in \FF_d^+}\). For any \(n \in \NN\) and \(Z \in \partial \BB^d_\nc\), we have that 
    \[
        \p_{\s(n)}^\#(Z^*)^* \p_{\s(n)}^\#(Z^*) = \sum_{\lvert \s \rvert = n} \p_\s(Z^*)^*\p_\s(Z^*)
    \]
    and this quantity is non-singular.
\end{lemma}  

\begin{proof}
    Let \(Z \in \partial \BB^d_\nc\); then \(I - Z_1Z_1^* + \cdots + Z_dZ_d^* = 0\), so the same manipulations as in \Cref{lem:ZerosOutside} and \Cref{lem:ZerosInside} this time yield the stated equality immediately. It remains to show that this quantity is non-singular on \(\partial \BB_\nc^d\).
    
    To see this, suppose to the contrary that
    \[
        \det \left(\p_{\s(n)}^\#(W^*)^*\p_{\s(n)}^\#(W^*)\right) = 0
    \]
    for some \(W \in \partial \BB^d_\nc\). Since for any matrix \(A\) we have \(\det AA^* = (\det A) (\overline{\det A}) = \lvert \det A \rvert^2\), this then implies that
    \[
        \det \p_{\s(n)}^\#(W^*) = 0,
    \]
    so in particular if \(W\) is \(m \times m\) then there exists some \(\xi \in \CC^{m}\setminus\{0\}\) such that \(\p_{\s(n)}^\#(W^*)\xi = 0\).

    Recall from the proof of \Cref{lem:ZerosOutside} that if \(Z \in \BB^d_\nc\) then
    \begin{multline*}
        \p_{\s(n)}^\#(Z^*)^* \p_{\s(n)}^\#(Z^*) = \sum_{\lvert \s \rvert = n} \p_\s(Z^*)^* \p_\s(Z^*) + (I - Z_1Z_1^* - \cdots - Z_dZ_d^*)  \\ + \sum_{1 \leq \lvert \s \rvert \leq n-1} \p_\s(Z^*)^*(I - Z_1Z_2^* - \cdots - Z_dZ_d^*)\p_\s(Z^*),
    \end{multline*}
    so that
    \[
        \p_{\s(n)}^\#(Z^*)^* \p_{\s(n)}^\#(Z^*) \geq I - \sum_{k=1}^d Z_kZ_k^*.
    \]
    Notice that if \(r \in (0,1)\) then \(rW \in \BB^d_\nc\), providing us with the matrix inequality
    \begin{align*}
        \p_{\s(n)}^\#(rW^*)^* \p_{\s(n)}^\#(rW^*) & \geq I - \sum_{k=1}^d (rW_k)(rW_k)^* \\
        & = \left(I - \left(\sum_{k=1}^d (rW_k)(rW_k)^*\right)^{1/2}\right)\left(I + \left(\sum_{k=1}^d (rW_k)(rW_k)^*\right)^{1/2}\right) \\
        & = \left(I - \left(r^2\sum_{k=1}^d W_kW_k^*\right)^{1/2}\right)\left(I + \left(r^2\sum_{k=1}^d W_kW_k^*\right)^{1/2}\right);
    \end{align*}
    recalling that \(W \in \partial \BB^d_\nc\), this becomes
    \[
        \p_{\s(n)}^\#(rW^*)^* \p_{\s(n)}^\#(rW^*) \geq (I - rI^{\frac12})(I + rI^{\frac12}) = (1-r)(1+r)I.
    \]
Thus, we have
    \[
        \frac{\xi^* \p_{\s(n)}^\#(rW^*)^* \p_{\s(n)}^\#(rW^*) \xi}{1-r} \geq (1+r)\xi^* \xi > \xi^* \xi > 0.
    \]
    Write the polynomial \(\p_{\s(n)}^\#\) as \(\sum_{\s \in \FF_d^+} \a_\s Z^\s\); with this notation, \(\p_{\s(n)}^\#(W^*)\xi = 0\) implies that 
    \[
        \a_\emptyset \xi = - \sum_{\s \neq \emptyset} \a_\s (W^*)^\s \xi.
    \]
    Substituting this in for the constant coefficient in \(\p_{\s(n)}^\#(rW^*)\xi\) we obtain
    \[
        \p_{\s(n)}^\#(rW^*)\xi = - \sum_{\s \neq \emptyset} \a_\s (W^*)^\s \xi + \sum_{\s \neq \emptyset} \a_\s (rW^*)^\s \xi = \sum_{\s\neq\emptyset} \a_\s (r^{\lvert \s \rvert} - 1) (W^*)^{\s} \xi,
    \]
    which in turns transforms our inequality into
    \begin{align*}
        0 & < \xi^* \xi < \frac{\left(\sum_{\s\neq\emptyset} \a_\s (r^{\lvert \s \rvert} - 1) (W^*)^{\s} \xi\right)^*\left(\sum_{\s\neq\emptyset} \a_\s (r^{\lvert \s \rvert} - 1) (W^*)^{\s} \xi\right)}{1-r} \\
        & = \frac{1}{1-r}\left((1-r)\sum_{\s\neq\emptyset} \a_\s \left(\sum_{j=0}^{\lvert \s \rvert - 1} r^j \right) (W^*)^{\s} \xi\right)^*\left((1-r)\sum_{\s\neq\emptyset} \a_\s \left(\sum_{j=0}^{\lvert \s \rvert - 1} r^j \right) (W^*)^{\s} \xi\right) \\
        & = (1-r)\left(\sum_{\s\neq\emptyset} \a_\s \left(\sum_{j=0}^{\lvert \s \rvert - 1} r^j \right) (W^*)^{\s} \xi\right)^*\left(\sum_{\s\neq\emptyset} \a_\s \left(\sum_{j=0}^{\lvert \s \rvert - 1} r^j \right) (W^*)^{\s} \xi\right) \\
        & \to 0 < \xi^*\xi
    \end{align*}
    as \(r \to 1\), a contradiction. Hence no such \(W\) can exist.
\end{proof}

\end{document}
